% ===== Preâmbulo compartilhado (Tectonic / XeTeX) — documentos da FTL094 =====
\documentclass[11pt,a4paper]{article}
\usepackage{fontspec}
\usepackage[brazil]{babel}
\usepackage{geometry}
\geometry{a4paper,margin=2.5cm}
\usepackage{amsmath,amssymb}
\usepackage{graphicx}
\usepackage{booktabs}
\usepackage{array}
\usepackage{longtable}
\usepackage{tabularx}
\usepackage{float}
\usepackage{enumitem}
\usepackage{xcolor}
\usepackage{tikz}
\usetikzlibrary{arrows.meta,positioning,automata,shapes.geometric,fit,backgrounds,calc}
\usepackage{pgfplots}
\pgfplotsset{compat=1.18}
\usepackage{listings}
\usepackage{caption}
\usepackage{fancyhdr}
\usepackage[hidelinks]{hyperref}

% ----- cores -----
\definecolor{ufamblue}{RGB}{0,51,102}
\definecolor{codegreen}{rgb}{0,0.5,0}
\definecolor{codeblue}{rgb}{0,0,0.7}
\definecolor{codered}{rgb}{0.6,0,0}
\definecolor{lgray}{rgb}{0.95,0.95,0.95}

% ----- listagens C -----
\lstdefinestyle{c}{language=C,basicstyle=\ttfamily\footnotesize,
  keywordstyle=\color{codeblue}\bfseries,commentstyle=\color{codegreen}\itshape,
  stringstyle=\color{codered},numbers=left,numberstyle=\tiny\color{gray},
  frame=single,breaklines=true,tabsize=2,showstringspaces=false,columns=fullflexible}

% ----- constantes do projeto -----
\newcommand{\projnome}{Sistema de Navegação Autônoma para o Robô Móvel Khepera~IV}
\newcommand{\disciplina}{FTL094 --- Engenharia de Software de Tempo Real}
\newcommand{\periodo}{Período Letivo 2026/1}
\newcommand{\versaodoc}{1.0}

% ----- cabeçalho/rodapé -----
\pagestyle{fancy}\fancyhf{}
\rhead{\small \tipodoc}
\lhead{\small Khepera~IV --- FTL094}
\cfoot{\thepage}
\renewcommand{\headrulewidth}{0.4pt}

% ----- coluna X centralizada p/ tabularx -----
\newcolumntype{C}{>{\centering\arraybackslash}X}

% ----- capa padronizada -----
\newcommand{\fazcapa}{%
\begin{titlepage}\centering
{\large\bfseries UNIVERSIDADE FEDERAL DO AMAZONAS (UFAM)}\\[3pt]
{\normalsize Centro de P\&D em Tecnologia Eletrônica e da Informação --- CETELI}\\[2pt]
{\normalsize Programa de Pós-Graduação em Engenharia Elétrica --- PPGEE}\\[26pt]
{\large \disciplina}\\[2pt]
{\normalsize \periodo}\\[64pt]
\rule{\textwidth}{1.2pt}\\[12pt]
{\Huge\bfseries\color{ufamblue} \tipodoc}\\[14pt]
{\Large \projnome}\\[10pt]
\rule{\textwidth}{1.2pt}\\[48pt]
{\large\bfseries Equipe}\\[8pt]
{\normalsize Gabriel Henrique de Souza Nazaré}\\[3pt]
{\normalsize Luan Alexandre Ramos Barreto}\\[3pt]
{\normalsize Matheus Rocha Canto}\\[3pt]
{\normalsize Matheus Serrão Uchôa}\\[3pt]
{\normalsize Ricardo Mendonça Braz}\\[3pt]
{\normalsize Yan Fernandes da Silva}\\[40pt]
{\normalsize Docente responsável: Prof. Dr.~[nome do docente]}\\
\vfill
{\small Documento de tipo \emph{\tipodoc} --- Versão \versaodoc}\\[2pt]
{\small Manaus --- AM, \datadoc}
\end{titlepage}}

% ----- tabela de controle de versões -----
\newcommand{\controleversoes}{%
\section*{Controle de versões}
\begin{tabularx}{\textwidth}{@{}l l l X@{}}
\toprule
\textbf{Versão} & \textbf{Data} & \textbf{Autor} & \textbf{Descrição} \\
\midrule
1.0 & \datadoc & Equipe FTL094 & Emissão inicial do documento. \\
\bottomrule
\end{tabularx}
\bigskip}

\newcommand{\tipodoc}{Análise de Requisitos}
\newcommand{\datadoc}{16 de junho de 2026}

\begin{document}
\fazcapa
\controleversoes
\tableofcontents
\newpage

% =====================================================================
\section{Introdução}
% =====================================================================
\subsection{Propósito}
Este documento especifica os requisitos do \emph{\projnome}. Destina-se à equipe de desenvolvimento, ao cliente (coordenação da disciplina) e aos avaliadores, servindo de base contratual para projeto, implementação e verificação.

\subsection{Escopo}
O produto é o software embarcado que dota o robô Khepera~IV de \textbf{navegação autônoma ponto-a-ponto com desvio de obstáculos}. O software lê sensores (proximidade e localização), decide a ação e comanda os motores, em laço de controle de tempo real, conduzindo o robô de uma posição inicial ao alvo sem colisões.

\subsection{Definições, acrônimos e abreviações}
\begin{description}[style=nextline,leftmargin=1.2cm,font=\ttfamily]
  \item[RF] Requisito funcional.
  \item[RNF] Requisito não funcional.
  \item[RT] Requisito temporal (tempo real).
  \item[UC] Caso de uso.
  \item[Pose] Posição $(x,y)$ e orientação ($\theta$) do robô.
\end{description}

\subsection{Referências}
IEEE Std 830-1998 (\emph{Recommended Practice for Software Requirements Specifications}); Sommerville (2019), cap.~4; Documento de Plano de Projeto (companheiro).

\subsection{Visão geral}
A Seção~\ref{sec:geral} descreve o produto em alto nível; a Seção~\ref{sec:esp} detalha os requisitos; a Seção~\ref{sec:modelos} traz casos de uso; a Seção~\ref{sec:rastr} apresenta a rastreabilidade.

% =====================================================================
\section{Descrição geral}\label{sec:geral}
% =====================================================================
\subsection{Perspectiva do produto}
O software opera embarcado no robô, interagindo com sensores, atuadores e o ambiente, conforme o diagrama de contexto da Figura~\ref{fig:ctx}.

\begin{figure}[H]\centering
\begin{tikzpicture}[
  box/.style={draw,rounded corners,minimum height=0.9cm,minimum width=2.6cm,align=center,font=\small},
  sys/.style={draw,thick,fill=ufamblue!10,rounded corners,minimum height=1.4cm,minimum width=3.6cm,align=center,font=\small\bfseries},
  >={Stealth}]
  \node[sys] (s) {Software de\\Navegação};
  \node[box] (sens) [left=2.4cm of s] {Sensores\\(IR, GPS, IMU)};
  \node[box] (act) [right=2.4cm of s] {Atuadores\\(motores das rodas)};
  \node[box] (op) [above=1.3cm of s] {Operador /\\Cliente};
  \node[box] (env) [below=1.3cm of s] {Ambiente\\(obstáculos, alvo)};
  \draw[->] (sens) -- node[above,font=\scriptsize]{leituras} (s);
  \draw[->] (s) -- node[above,font=\scriptsize]{velocidades} (act);
  \draw[->] (op) -- node[right,font=\scriptsize]{alvo, parâmetros} (s);
  \draw[->] (s) -- node[right,font=\scriptsize]{telemetria} (op);
  \draw[->] (env) -- node[right,font=\scriptsize]{estímulos} (sens.south);
  \draw[->] (act.south) |- node[below,font=\scriptsize]{movimento} (env.east);
\end{tikzpicture}
\caption{Diagrama de contexto do sistema.}
\label{fig:ctx}
\end{figure}

\subsection{Funções do produto}
Em alto nível, o produto: (i) estima a pose do robô; (ii) calcula a direção ao alvo; (iii) detecta obstáculos por proximidade; (iv) decide entre seguir ao alvo, girar ou contornar; (v) comanda as rodas; (vi) detecta a chegada e para; e (vii) registra telemetria.

\subsection{Características dos usuários}
O usuário-operador possui conhecimento técnico básico de robótica; interage definindo o alvo e os parâmetros e observando a telemetria. Não há interação durante a missão (operação autônoma).

\subsection{Restrições gerais}
Linguagem C; plataforma Khepera~IV; sensoriamento de curto alcance; ausência de mapa prévio; execução em laço periódico de tempo real.

\subsection{Suposições e dependências}
Ambiente estático e plano; alvo alcançável (existe caminho livre); disponibilidade da localização (GPS/IMU em simulação; odometria no robô real).

% =====================================================================
\section{Requisitos específicos}\label{sec:esp}
% =====================================================================
\subsection{Requisitos funcionais}
\begin{table}[H]\centering
\caption{Requisitos funcionais.}
\renewcommand{\arraystretch}{1.2}
\begin{tabularx}{\textwidth}{@{}l X c@{}}
\toprule
\textbf{Id} & \textbf{Descrição} & \textbf{Prior.} \\
\midrule
RF-01 & O sistema deve conduzir o robô de uma posição inicial a uma posição-alvo definida por coordenadas. & Essencial \\
RF-02 & O sistema deve desviar autonomamente de obstáculos estáticos presentes no percurso. & Essencial \\
RF-03 & O sistema deve detectar a chegada ao alvo (dentro de tolerância) e parar com segurança. & Essencial \\
RF-04 & O sistema deve estimar a pose do robô (posição e orientação). & Essencial \\
RF-05 & O sistema deve comandar as velocidades das rodas (tração diferencial). & Essencial \\
RF-06 & O sistema deve recuperar-se de situações de aprisionamento (mínimo local). & Importante \\
RF-07 & O sistema deve registrar telemetria periódica (estado, pose, sensores, comandos). & Importante \\
RF-08 & O sistema deve permitir a configuração do alvo e dos parâmetros de controle. & Desejável \\
RF-09 & O sistema deve sinalizar/identificar o estado interno de navegação. & Desejável \\
\bottomrule
\end{tabularx}
\end{table}

\subsection{Requisitos não funcionais}
\begin{table}[H]\centering
\caption{Requisitos não funcionais.}
\renewcommand{\arraystretch}{1.2}
\begin{tabularx}{\textwidth}{@{}l l X@{}}
\toprule
\textbf{Id} & \textbf{Categoria} & \textbf{Descrição} \\
\midrule
RNF-01 & Segurança & Operar sem colisões com obstáculos ou paredes. \\
RNF-02 & Desempenho & Alcançar o alvo em tempo finito e limitado para o cenário especificado. \\
RNF-03 & Confiabilidade & Comportamento determinístico e repetível entre execuções. \\
RNF-04 & Portabilidade & Permitir migração simulação $\to$ robô físico, isolando a localização. \\
RNF-05 & Manutenibilidade & Código modular, com parâmetros centralizados e documentados. \\
RNF-06 & Observabilidade & Prover telemetria persistente para diagnóstico (registro em arquivo). \\
\bottomrule
\end{tabularx}
\end{table}

\subsection{Requisitos temporais (tempo real)}
\begin{table}[H]\centering
\caption{Requisitos de tempo real.}
\renewcommand{\arraystretch}{1.2}
\begin{tabularx}{\textwidth}{@{}l X@{}}
\toprule
\textbf{Id} & \textbf{Descrição} \\
\midrule
RT-01 & O laço de controle deve executar como tarefa periódica de período $T=16$~ms. \\
RT-02 & A cada ciclo, a sequência percepção$\to$decisão$\to$atuação deve respeitar o \emph{deadline} $D=T$. \\
RT-03 & O WCET do ciclo deve ser significativamente inferior a $T$, garantindo folga para \emph{jitter}. \\
\bottomrule
\end{tabularx}
\end{table}

\subsection{Requisitos de interface}
\textbf{Hardware:} 5 sensores IR frontais/laterais (proximidade); GPS e IMU (localização, em simulação); dois motores de roda em modo de velocidade. \textbf{Software:} API de controladores do Webots (\texttt{wb\_*}); arquivo de telemetria em texto.

% =====================================================================
\section{Modelos: casos de uso}\label{sec:modelos}
% =====================================================================
\begin{figure}[H]\centering
\begin{tikzpicture}[uc/.style={draw,ellipse,minimum width=3.2cm,minimum height=1cm,align=center,font=\small},>={Stealth}]
  % ator
  \node (a) at (0,0) {\begin{tikzpicture}[scale=0.5]
    \draw[thick] (0,1) circle(0.25);
    \draw[thick] (0,0.75)--(0,-0.1);
    \draw[thick] (-0.4,0.5)--(0.4,0.5);
    \draw[thick] (0,-0.1)--(-0.35,-0.7);
    \draw[thick] (0,-0.1)--(0.35,-0.7);
  \end{tikzpicture}};
  \node[below=0cm of a,font=\small] {Operador};
  % fronteira do sistema
  \node[uc] (u1) at (4,1.6) {UC-01\\Executar missão};
  \node[uc] (u2) at (4,0) {UC-02\\Configurar alvo/parâm.};
  \node[uc] (u3) at (4,-1.6) {UC-03\\Monitorar telemetria};
  \begin{scope}[on background layer]
    \node[draw,dashed,fit=(u1)(u2)(u3),inner sep=10pt,label=above:{\small Sistema de Navegação}] {};
  \end{scope}
  \draw (a) -- (u1); \draw (a) -- (u2); \draw (a) -- (u3);
\end{tikzpicture}
\caption{Diagrama de casos de uso.}
\label{fig:uc}
\end{figure}

\noindent\textbf{UC-01 --- Executar missão de navegação.}
\emph{Ator:} Operador. \emph{Pré-condição:} alvo configurado; robô na posição inicial. \emph{Fluxo principal:} (1) operador inicia a execução; (2) o sistema entra no laço de controle; (3) a cada ciclo estima a pose, avalia obstáculos e atua; (4) ao detectar a chegada, para. \emph{Pós-condição:} robô no alvo, sem colisões. \emph{Fluxo alternativo:} obstáculo frontal $\Rightarrow$ contorno; mínimo local $\Rightarrow$ manobra de escape.

\noindent\textbf{UC-02 --- Configurar alvo e parâmetros.} O operador define coordenadas do alvo e ganhos/limiares antes da missão.

\noindent\textbf{UC-03 --- Monitorar telemetria.} O operador acompanha o registro periódico (estado, pose, distância, sensores).

% =====================================================================
\section{Matriz de rastreabilidade}\label{sec:rastr}
% =====================================================================
A Tabela~\ref{tab:rastr} relaciona requisitos, objetivos e casos de teste (definidos no Plano de Testes), assegurando cobertura e verificabilidade.

\begin{table}[H]\centering
\caption{Rastreabilidade requisitos $\leftrightarrow$ casos de uso $\leftrightarrow$ testes.}
\label{tab:rastr}
\begin{tabularx}{\textwidth}{@{}l X l l@{}}
\toprule
\textbf{Requisito} & \textbf{Verificação} & \textbf{UC} & \textbf{Teste} \\
\midrule
RF-01 & Robô atinge o alvo em cenário livre & UC-01 & CT-01 \\
RF-02 & Desvio de obstáculo único e múltiplo & UC-01 & CT-02, CT-03 \\
RF-03 & Parada dentro da tolerância & UC-01 & CT-01 \\
RF-06 & Escape de alvo atrás de parede & UC-01 & CT-04 \\
RF-07 & Telemetria gravada em arquivo & UC-03 & CT-06 \\
RNF-01 & Ausência de colisões no slalom & UC-01 & CT-03 \\
RT-01 & Período de laço medido $=16$~ms & --- & CT-05 \\
\bottomrule
\end{tabularx}
\end{table}

\section{Critérios de aceitação}
O produto será aceito se: (i) atingir o alvo no cenário de 5 obstáculos sem colisões (RF-01, RF-02, RNF-01); (ii) parar dentro da tolerância (RF-03); (iii) respeitar o período de tempo real (RT-01); e (iv) registrar telemetria verificável (RF-07).

\end{document}
